\chapter{Grading-совместимое endpoint-расширение семейной тройки}
\label{ch:v8-baryon-c0-grading-compatible-family-triplet-endpoint-extension-gate}

Текущий charged-singlet target имеет разложение по grading
\begin{equation}
 H_{e_R}=H_-oplus H_+,
 \qquad \dim H_-=1,\qquad \dim H_+=2,
 \qquad \Gamma_{e_R}=(-I_1)\oplus(+I_2).
 \label{eq:v8-baryon-c0-current-target-grading-split}
\end{equation}
Неприводимая вещественная семейная тройка допускает только скалярную
градуировку:
\begin{equation}
 [\Gamma_{mathbf3},J_a]=0\ \forall a
 \quad\Longrightarrow\quad
 \Gamma_{mathbf3}=\varepsilon I_3,
 \qquad \varepsilon\in\{+1,-1\}.
 \label{eq:v8-baryon-c0-family-triplet-scalar-grading}
\end{equation}
Поэтому числа новых target-состояний в двух ветвях равны
\begin{equation}
 n_+=3-\dim H_+=1,
 \qquad n_-=3-\dim H_-=2.
 \label{eq:v8-baryon-c0-homogeneous-triplet-extension-counts}
\end{equation}

\section{Единственная минимальная ветвь}

В положительной ветви добавляется одна линия того же
charged-singlet gauge-типа:
\begin{equation}
 H_{24}=H_{23}\oplus\mathbb C e_R^{(t,0)},
 \qquad \Gamma e_R^{(t,0)}=+e_R^{(t,0)},
 \qquad Y e_R^{(t,0)}=-e_R^{(t,0)}.
 \label{eq:v8-baryon-c0-minimal-plus-target-extension}
\end{equation}
Она завершает target-тройку
\begin{equation}
 H_{mathbf3}^{(+)}=operatorname{span}_{\mathbb C}
 \{e_R^{(t,0)},e_R^{(t,1)},e_R^{(t,2)}\},
 \qquad \Gamma\big|_{H_{mathbf3}^{(+)}}=+I_3.
 \label{eq:v8-baryon-c0-plus-family-target-triplet}
\end{equation}
Соответствующий connector-triplet использует единственную source-линию
$a_0$:
\begin{equation}
 W_{mathbf3}^{(+)}=operatorname{span}_{\mathbb C}
 \{|e_R^{(t,k)}\rangle\langle a_0|:k=0,1,2\},
 \qquad \dim_{\mathbb C}W_{mathbf3}^{(+)}=3.
 \label{eq:v8-baryon-c0-plus-family-connector-triplet}
\end{equation}
Все три стрелки нечётны, поскольку
\begin{equation}
 (+I_3)V+V(-1)=0,
 \qquad V\in W_{mathbf3}^{(+)}.
 \label{eq:v8-baryon-c0-plus-triplet-oddness}
\end{equation}
Стандартное семейное действие одновременно совместимо с grading,
гиперзарядом и Real-структурой:
\begin{equation}
 [J_a,+I_3]=[J_a,-I_3]=0,
 \qquad \overline{J_a}=J_a.
 \label{eq:v8-baryon-c0-plus-triplet-structural-compatibility}
\end{equation}
Лемма Шура снова даёт единственную изометрическую карту с точностью до
нефизического общего знака:
\begin{equation}
 \operatorname{Hom}_{SO(3)}(\mathbf3_{\rm fam},W_{mathbf3}^{(+)})
 =\mathbb RI_3,
 \qquad T^{\mathsf T}T=I_3\Longrightarrow T=\pm I_3.
 \label{eq:v8-baryon-c0-plus-triplet-schur-map}
\end{equation}

Отрицательная ветвь также формально допустима, но требует двух линий:
\begin{equation}
 H_{25}=H_{23}\oplus\mathbb C e_R^{(s,1)}
 \oplus\mathbb C e_R^{(s,2)},
 \qquad \Gamma\big|_{H_{mathbf3}^{(-)}}=-I_3,
 \qquad W_{mathbf3}^{(-)}=H_{mathbf3}^{(-)}\langle s_0|.
 \label{eq:v8-baryon-c0-minus-family-triplet-branch}
\end{equation}
При сохранении всех старых grading-знаков и gauge-типов имеем строгую
минимальность
\begin{equation}
 n_+=1<n_-=2,
 \qquad \text{поэтому минимальная ветвь единственна и выбирает }a_0.
 \label{eq:v8-baryon-c0-minimal-branch-selects-a0}
\end{equation}

\section{Цена endpoint-ковариантности}

Семейное действие не сохраняет три координатных endpoint-проектора, но
сохраняет их сумму:
\begin{equation}
 [J_a,P_k]\ne0\ \text{для некоторых }(a,k),
 \qquad [J_a,P_{mathbf3}]=0,
 \qquad P_{mathbf3}=P_0+P_1+P_2=I_3.
 \label{eq:v8-baryon-c0-family-triplet-endpoint-projectors}
\end{equation}
Ковариантное ассоциативное замыкание диагональной endpoint-алгебры и
семейных генераторов является полным блоком:
\begin{equation}
 \operatorname{Alg}(D_3,J_1,J_2,J_3)=M_3(\mathbb C),
 \qquad \dim_{\mathbb C}M_3=9,
 \qquad \Delta N_{\rm Herm}=6.
 \label{eq:v8-baryon-c0-family-triplet-endpoint-m3-closure}
\end{equation}
Следовательно, новый state недостаточен сам по себе: требуется также
семейно-ковариантный endpoint-блок.

Старый коннектор $|e_R^{(s)}\rangle\langle s_0|$ остаётся семейным
синглетом, поэтому полный connector-носитель имеет тип
\begin{equation}
 W_{\rm conn}^{\rm ext}\simeq\mathbf1\oplus\mathbf3.
 \label{eq:v8-baryon-c0-extended-connector-one-plus-three}
\end{equation}
Его инвариантная симметричная covariance содержит два независимых веса:
\begin{equation}
 C_{SO(3)}=\gamma_1P_{mathbf1}+\gamma_3P_{mathbf3},
 \qquad \dim\{C=C^{\mathsf T}:[C,SO(3)]=0\}=2.
 \label{eq:v8-baryon-c0-singlet-triplet-covariance-two-weights}
\end{equation}
Значит, условная тройка выбирает карту внутри triplet, но не отношение её
скорости к оставшемуся синглету.

Итоговый extension ledger равен
\begin{equation}
 N_{\rm extension}=0/3:
 \quad\{e_R^{(t,0)},\ M_3\text{-endpoint},\ SO(3)\text{-covariant frame}\}.
 \label{eq:v8-baryon-c0-family-triplet-extension-ledger}
\end{equation}

\begin{theorem}[Минимальное grading-совместимое завершение]
При сохранении grading и gauge-типов всех старых состояний единственное
минимальное завершение семейной тройки добавляет одну положительно
градуированную charged-singlet линию к двум существующим линиям $H_+$.
Оно выбирает source $a_0$ и условно даёт интертвейнер $T=\pm I_3$.
\end{theorem}

\begin{nogocriterion}[Минимальность расширения не есть parent-origin]
Счёт $1<2$ выбирает наименьшую архитектуру, но не создаёт новое состояние,
полный endpoint-блок $M_3$ или его ковариантный шумовой кадр. До их вывода
из одного действия физическая карта $c_0$ не считается полученной.
\end{nogocriterion}

\begin{protocol}\begin{enumerate}
\item текущие grading-кратности: \textbf{$1_-+2_+$};
\item новых states в плюс-ветви: \textbf{$1$};
\item новых states в минус-ветви: \textbf{$2$};
\item минимальная ветвь: \textbf{плюс};
\item выбранный source: \textbf{$a_0$};
\item расширенный carrier: \textbf{$H_{24}$};
\item условный family-Hom: \textbf{$\mathbb RI_3$};
\item endpoint-замыкание: \textbf{$M_3(\mathbb C)$};
\item полный connector-тип: \textbf{$\mathbf1\oplus\mathbf3$};
\item независимых invariant rates: \textbf{$2$};
\item выведенных структур: \textbf{$0/3$};
\item следующий гейт: \textbf{селектор относительной скорости singlet--triplet}.
\end{enumerate}\end{protocol}

Машинный сертификат сохранён в
\path{s2t_v8_baryon_c0_grading_compatible_family_triplet_endpoint_extension_gate_results.json}.